\scanpage{23}
\deleted{\uline{Lemme}} \inserted{\uline{Proposition 1}} $S$ préschéma, $F:(\mathrm{Sch})_{/S}^{\circ}\mathbin{-}\mathrm{Ens}$, $X$ \deleted{$\in\operatorname{Ob}(\mathrm{Sch})_{/S}$} \inserted{un préschéma loc. de prés. finie sur $S$}, $\xi\in F(X)$. On suppose

$\underset{\leftarrow}{L}_1)$ \deleted{\unsure{a)}} $F$ de nature locale

$\underset{\leftarrow}{L}_2)$ \deleted{\unsure{b)}} $F$ ``commute aux limites inductives d'anneaux''

$R_{\mathrm{loc}})$ \deleted{\unsure{c)}} $(X,\xi)$ représente $F$ quand on se restreint aux arguments $T/S$ qui sont \deleted{des anneaux} \inserted{locaux} et essentiellement de type fini sur $S$, à extension résiduelle finie sur $S$.
Alors $(X,\xi)$ représente $F$, i.e. $\xi:X\mathbin{-}F$ est un isom.

\margin{Considérer \unsure{le} cas $F\to G$ qui doit être \uline{mono}, \uline{imm. ouverte}, \uline{isom.}}

\uline{Dém} Comme $F$ de nature locale, on peut supposer $S$ affine (aff), alors $R_{\mathrm{loc}}$ signifie que pour $\forall$ \inserted{spectre $T$ d'un} localisé $B$ d'une algèbre de t.f. sur $A$,
\[
\operatorname{Hom}_S(\deleted{\operatorname{Spec}B}\inserted{T},X)\simeq F(T)
\]
par $\xi$. \inserted{un isom} On est ramené à prouver que ceci reste vrai pour $\forall T$ de t.f. sur $S$ (aff).

1°) \uline{Injectivité} Soient $u,v:T\rightrightarrows X$ \deleted{\unsure{de f}} tels que $F(u)(\xi)=F(v)(\xi)$, prouvons $u=v$. Par l'hyp. $R_{\mathrm{loc}}$ $u$ et $v$ sont $=$ restreints aux $\operatorname{Spec}\mathcal O_{T,t}$, d'où la conclusion.

2°) Surjectivité. Soit $\unsure{\eta|y}\in F(T)$. Pour $\forall t\in T$ \inserted{à ext. résiduelle finie}, il existe grâce à $R_{\mathrm{loc}}$ un $u_t:\operatorname{Spec}\mathcal O_{T,t}\to X$ (unique) induisant $\xi_t\in F(\operatorname{Spec}\mathcal O_{T,t})$. \deleted{En outre}
\mspage{22}
